\documentclass[11pt]{amsart}
\usepackage[margin=1.04in]{geometry}
\usepackage[T1]{fontenc}
\usepackage{lmodern,amsmath,amssymb,amsfonts,mathtools,mathrsfs,booktabs,longtable,array,microtype}
\usepackage[hidelinks]{hyperref}
\ifdefined\pdfinfoomitdate\pdfinfoomitdate=1\fi
\ifdefined\pdftrailerid\pdftrailerid{}\fi
\ifdefined\pdfsuppressptexinfo\pdfsuppressptexinfo=15\fi
\newtheorem{theorem}{Theorem}[section]
\newtheorem{proposition}[theorem]{Proposition}
\newtheorem{lemma}[theorem]{Lemma}
\newtheorem{corollary}[theorem]{Corollary}
\theoremstyle{definition}\newtheorem{definition}[theorem]{Definition}
\newtheorem{example}[theorem]{Example}
\theoremstyle{remark}\newtheorem{remark}[theorem]{Remark}
\newcommand{\E}{\mathbb E}\newcommand{\Pp}{\mathbb P}\newcommand{\R}{\mathbb R}
\newcommand{\N}{\mathbb N}\newcommand{\one}{\mathbf 1}
\newcommand{\tr}{\operatorname{tr}}\newcommand{\TV}{\operatorname{TV}}
\newcommand{\conv}{\operatorname{conv}}\newcommand{\supp}{\operatorname{supp}}
\newcommand{\Law}{\operatorname{Law}}\newcommand{\rank}{\operatorname{rank}}
\newcommand{\diag}{\operatorname{diag}}\newcommand{\osc}{\operatorname{osc}}
\newcommand{\Lip}{\operatorname{Lip}}\newcommand{\dist}{\operatorname{dist}}
\newcommand{\calR}{\mathcal R}\newcommand{\calE}{\mathfrak E}
\newcommand{\norm}[1]{\left\lVert#1\right\rVert}
\newcommand{\abs}[1]{\left|#1\right|}\newcommand{\HS}{\mathrm{HS}}
\DeclareMathOperator*{\esssup}{ess\,sup}
\title[General Theta Foundations I]{General Theta Foundations I:\\Acquired Geometry and Causal Resource Transfer}
\author{Qian Qi}\date{11 October 2026}
\subjclass[2020]{Primary 62B15, 93E20; Secondary 94A34, 60J05, 81P45}
\keywords{causal experiment, acquired geometry, controlled observation, predictive quotient, finite-state control, quantum instruments}
\begin{document}\raggedbottom
\begin{abstract}
We study known controlled causal experiments with a finite executable affine test closure and a bounded terminal predictive task. A feedback-uniform overlap of the actually acquired predictive laws, together with an interior preparation buffer, yields decaying Bellman sensitivity without a one-step Wasserstein contraction. Combining this mechanism with conditional calibration and an acquired-mass curvature estimate gives a matching finite-resource theorem against the full-history adaptive optimum. For $n$ paid reports, checkpoint excess risk is comparable to $m^{-2/d}$; a machine with $M$ total states that enforces its own exact deadline has excess comparable to $(M-n)^{-2/d}$, uniformly in $n$ and $M$. All control, phase, stopping and readout states are counted. The converse covers every legal randomized Borel controller. Active Gaussian sensors with arbitrary positive invertible hidden-state transition, and noncommuting quantum instruments with feedback unitaries, verify the same theorem from their raw kernels. The Gaussian class includes an explicitly Wasserstein-expansive report step; both classes exhibit strict observation-dependent action switching. Visible rank-changing acquisition yields a weighted singular resource profile. We also prove a matched same-task calibration--deficiency extension with a protected post-control bit, and distinguish finite-state existence from certified finite-precision execution.
\end{abstract}\maketitle
\section{Introduction}\label{sec:intro}
An experiment does not supply a geometric state space to its observer. It supplies preparations, admissible interventions and reports. A state is predictive only to the extent that it determines executable future tests; a geometric resolution is useful only to the extent that it is acquired with positive probability and can be retained by a causal machine. The problem is to pass from these primitive data to a sharp resource law without treating a formal reachable set, a free exact posterior, or a phase-dependent program as an acquired finite resource.

There are two distinct obstacles. Static quantization may discard information needed to choose subsequent observations. Even when a compressed state predicts the next report well, the continuation value of choosing the wrong action can be nonsmooth. A pointwise Lipschitz estimate then loses an order relative to the quadratic terminal task. Moreover, stability of a passive filter does not by itself compare a finite-state feedback controller with the optimal full-history controller of the same experiment.

The main theorem resolves these obstacles on a specified, nontrivial class of controlled positive instruments. Its new stability mechanism is an overlap of the \emph{actual acquired laws}, uniform over feedback choices. If every controlled block retains a common positive component, the oscillation of every terminal Bellman continuation decays. Concavity and an interior preparation buffer convert this oscillation bound into a local Lipschitz bound. An acquired-density estimate then integrates the nonsmooth supporting-plane error to second order. The resulting controller chooses a Bellman-optimal action at its own calibrated retained state, not at an inaccessible full-history posterior.

Write $V_0(u_0)$ for the unrestricted full-history adaptive optimum of the declared terminal task. Under the hypotheses of Theorem~\ref{thm:main}, for every $n\ge1$,
\begin{equation}\label{eq:intro-law}
 \calR_{\rm cp}(n,m)-V_0(u_0)\asymp m^{-2/d},\qquad
 \calR_{\rm aut}(n,M)-V_0(u_0)\asymp(M-n)^{-2/d}.
\end{equation}
The second class is nonempty precisely when $M\ge n+1$. Its $n$-state charge is an internal exact-deadline charge, not a change from labels to bits. The constants are uniform in both budgets for fixed experiment certificates. They need not remain bounded as acquired overlap disappears. The baseline concerns the same dynamic terminal experiment for each $n$; no monotone static-parameter learning curve is inferred.

Two raw-kernel consequences give the theorem operational content. In the classical experiment, an observer selects among Gaussian sensors before each report, while the hidden state subsequently evolves through any fixed positive invertible transition. The theorem includes sensor steps whose acquired belief kernel expands Euclidean Wasserstein distance. In the quantum experiment, a report-dependent posterior is followed by a chosen unitary and a strictly positive channel. Small-effect report corridors yield feedback-uniform overlap even when one-step transport estimates do not contract. An asymmetric informationally complete detector makes the feedback choice strictly depend on the current acquired state. These results concern known models, finite action sets, and a finite terminal task; they do not assert unknown-kernel learning, arbitrary stopping, or universal control optimality beyond this interface.

The conditional-calibration identity, its Bellman--Jensen form, the integrated convex-curvature estimate, and the internal-clock argument have antecedents in the retained work \cite{GTFControl,GTFAllocation}. The likelihood-normalized three-defect identity was established separately in \cite{GTFBayes}. We reprove the needed results here so that the native article is self-contained. The new connection is the route from feedback-uniform \emph{acquired overlap} to the sharp controlled total-state law, together with its active Gaussian, feedback quantum and rank-changing verifications. Neither a conditional projection identity nor the classical common-mass contraction principle is claimed as a new general fact.

The distinction from value-directed belief compression \cite{PB2000,PB2001,PB2002}, finite-memory feedback approximation \cite{KY2022}, and recent belief-window and partition-filter approximations \cite{Kim2026,DGM2026} is the joint guarantee: an implemented autonomous controller is compared with the same full-history optimum, while a converse applies to every randomized Borel controller with the stated total alphabet. Section~\ref{sec:comparison} makes the comparison at the level of tasks, memory objects and theorem hypotheses. Section~\ref{sec:resources} treats numerical execution and causal simulation separately from hard-state existence.

\section{Experiments, predictive realization and admissible resources}\label{sec:experiments}
\subsection{Primitive data and predictive quotient}
All quotient assertions are made at a fixed experiment interface phase. In a time-inhomogeneous experiment the phase belongs to that interface; an autonomous controller must still carry its own phase through its charged labels.

A prepared causal experiment consists of standard Borel report spaces $Y_a$, a finite nonempty action set $A$, positive causal probability kernels for successive reports and protected terminal tests, and a preparation kernel. A legal strategy selects the next action from the observed history and independent private randomness. A future test is a bounded function of a legally executable finite continuation, including its prescribed interventions and terminal readout. Two prepared histories are equivalent when all such future test expectations coincide. Convex preparation means an actual random mixture of preparations, not a formal interpolation of an unimplemented state.

We impose the following finite realization hypothesis. A finite collection of executable tests and the constant test spans a separating affine space of dimension $d+1$. Their expectation vector $x(h)$ is Borel. Its attained image is Borel, with a Borel realization section. The expectations of all executable future tests factor through this vector. The unnormalized joint expectation of any coordinate test and the next report lies in the same affine span. A compact convex set $D\subset\R^d$, with nonempty interior in these coordinates, consists of realizable mixed preparations; all these affine identities hold on $D$. Norms and volumes below refer to this affine test chart, not to an ambient physical parameterization.

Consequently, for each $a$ and phase $t$ there are a reference probability measure $\lambda_{t,a}$ and Borel affine functions
\begin{equation}\label{eq:raw}
 \ell_{t,a}(y,u)=\ell^0_{t,a}(y)+c_{t,a}(y)u,\qquad
 U_{t,a}(y,u)=b_{t,a}(y)+A_{t,a}(y)u
\end{equation}
such that the report law is $\ell_{t,a}(y,u)\lambda_{t,a}(dy)$ and the normalized predictive update is
\begin{equation}\label{eq:posterior}
 F_{t,a}(y,u)=U_{t,a}(y,u)/\ell_{t,a}(y,u).
\end{equation}
Here $c_{t,a}(y)$ is a row linear functional and $c_{t,a}(y)v$ is scalar. For each $(t,a)$, the identities hold for all $u\in D$ off \emph{one} reference-null report set. When $\ell=0$, positivity and bounded coordinate tests force $U=0$ at the relevant state; the identity-level formulas use the closed perspective convention. The sharp theorem below uses positive common support. This hypothesis is verified from raw kernels in Sections~\ref{sec:classical} and~\ref{sec:quantum}; finite physical dimension alone would not prove it.

\begin{lemma}[Borel predictive realization and minimality]\label{lem:quotient}
Under the finite realization hypothesis, the attained test image, with its Borel structure, realizes the predictive quotient. Every other Borel statistic sufficient for all executable future tests determines this image by a Borel map on its realized image, provided it has a Borel realization section. The report and update kernels descend jointly in the finite action coordinate.
\end{lemma}
\begin{proof}
Equality of test coordinates implies equality of all future tests by separation and the factorization hypothesis; the converse follows because the coordinates themselves are executable. The Borel section identifies the quotient with the attained image, not with all of $D$. For another statistic $z$, compose its realization section with $h\mapsto x(h)$; sufficiency makes this independent of the representative. This gives the asserted Borel factorization. Formula~\eqref{eq:raw}, the common null set, and the finite action set give jointly Borel report and normalized update kernels. Iterated kernel integration reproduces every finite strategy and terminal test, proving realization. No unproved assertion that every measurable quotient is standard Borel is used.
\end{proof}

The larger $D$ is a domain of realizable mixtures used for uniform certificates. It is not asserted to be the actually acquired support. For an incoming $u$ define the acquired row law
\begin{equation}\label{eq:row}
 \mu_{t,a,u}=(F_{t,a}(\cdot,u))_\#
       (\ell_{t,a}(\cdot,u)\lambda_{t,a}).
\end{equation}
Under a full-history strategy $\pi$, the actual predictive vector $X_t$ and its law $\nu_t^\pi$ are obtained by iterating the chosen kernels. Under a finite controller, the relevant row mixture instead uses that controller's actual conditional barycenters. These two laws need not agree.

\subsection{The task and the two memory architectures}
The terminal audit is a bounded executable vector $Z$ satisfying
\begin{equation}\label{eq:task}
 \E[Z\mid\text{predictive state }u]=u,\qquad
 \E[\norm Z^2\mid u]=c(u),\qquad
 g(u)=c(u)-\norm u^2,
\end{equation}
where $c$ is affine. The response $r\in D$ is committed before the protected audit, and the loss is $\norm{Z-r}^2$. We also allow an audit in a larger Euclidean space with mean $i(u)$, where $i(u)=z_0+Ju$ is an affine isometry of the predictive chart. For loss $\norm{Z-i(r)}^2$, put $c(u)=\E[\norm{Z-z_0}^2\mid u]$. Direct expansion gives the same conditional loss $c(u)-2\langle u,r\rangle+\norm r^2$ and the same $g$. Thus audit noise orthogonal to the mean image is retained in the affine second moment, not discarded. This identifies the ordinary categorical or informationally complete POVM Brier task exactly with the displayed algebra; an alternative affine chart only changes fixed norm constants. All risks use this one physical task metric.

For fixed $n$, a checkpoint strategy may use the entire history to choose its $n$ acquisitions, then replace that history by at most $m$ labels before the audit. Its acquisition workspace is deliberately unrestricted and is not identified with the checkpoint label budget. Write its optimum as $\calR_{\rm cp}(n,m)$.

An autonomous controller has one finite set $S$ of at most $M$ labels, one initial label, a halt status and action kernel at every nonhalt label, a Borel next-label kernel depending only on the current label, chosen action and fresh report, and a terminal response attached to each halt label. All its phase and control information is in $S$. Fresh random coins are discarded after the update; a retained random seed would count as state. There is no external time, external last-report signal, re-readable history, persistent real register or free interim response register. The machine must halt after exactly $n$ reports almost surely. The underlying experiment supplies further reports if requested, so the horizon is not signaled by ending the physical input. Its optimum is $\calR_{\rm aut}(n,M)$; an empty class has value $+\infty$.

Known experiment constants and programmed conditional centroids are read-only program data, not free bit complexity. This definition concerns hard-state existence. Section~\ref{sec:resources} specifies additional precision, program, workspace and primitive-cost certificates for finite-bit execution. Preparation, each raw report, and the protected audit are separate physical calls: the core count is $N=n+2$.

The full-history optimum is given by
\begin{equation}\label{eq:bellman}
 V_n=g,\qquad V_{t-1}(u)=\min_{a\in A}\int V_t(w)\mu_{t,a,u}(dw).
\end{equation}
All full-history randomized strategies for the same preparation, action family and audit are admitted. The positivity and continuity conditions below ensure Borel minimizing selectors. A fixed strategy has an affine physical risk before optimization over its terminal response, so $V_t$ is concave. Thus~\eqref{eq:bellman} is a control optimum, not a same-controller substitution of approximate beliefs inside a cost.

\section{The controlled acquired-geometry theorem}\label{sec:main}
\begin{definition}[Acquired-overlap class]\label{def:class}
Fix the experiment and task of Section~\ref{sec:experiments}. The following certificates do not depend on $n$, a memory budget, or a selected controller.

\emph{Geometry and regularity.} There is a compact convex $K\subset\operatorname{int}D$ and $r>0$ with $K+B_r\subset D$, and $F_{t,a}(y,D)\subset K$ off the common null sets. Every acquired row~\eqref{eq:row}, for every $u\in D,t,a$, has density at most $\Lambda$ with respect to $d$-dimensional volume. The controlled row kernels are weakly continuous on $D$. The report laws obey
\[
 \TV(\ell_{t,a}(\cdot,u)\lambda_{t,a},
       \ell_{t,a}(\cdot,v)\lambda_{t,a})\le L_{\rm rep}\norm{u-v}.
\]
There are global Borel grids on $K$ of at most $m$ cells and diameter at most $C_Km^{-1/d}$, with legal representatives; compact convex $K$ supplies such grids.

\emph{Feedback-uniform acquired overlap.} For some integer $b\ge1$ and $\epsilon_*\in(0,1]$, and each block start $s$, there is a probability $\zeta_s$ on $K$ such that
\begin{equation}\label{eq:overlap}
 \Law_u^\pi(X_{s+b})\ \ge\ \epsilon_*\zeta_s
 \quad\text{for every }u\in D
 \text{ and every legal $b$-step history-feedback policy }\pi.
\end{equation}
The parameter $\epsilon_*$ is a mass certificate, not a simulator error. All intermediate reports and chosen actions enter the law on the left.

\emph{Autonomous support.} For each action, the report laws at every possible finite history generated by every legal program, at every phase, are equivalent to a fixed reference probability $\lambda_a$. Possible label paths have edges of positive reference-measure support. The prescribed report count does not appear as a free report mark.
\end{definition}

The density quantifier includes every retained barycenter a controller can create, not just the full-history trajectory. The overlap certificate concerns distributions actually obtained through reports, not a common subset of a formal reachable set. Neither density nor overlap is inferred from positivity alone. The common-support clause is needed only for the internally clocked conclusion. The exact identities below require neither density, overlap, nor a clock.

\begin{theorem}[Controlled acquired-geometry transfer]\label{thm:main}
Suppose Definition~\ref{def:class} and the finite realization and task hypotheses hold, and $u_0\in D$. For every integer $n\ge1$, $m\ge1$, and $M\ge n+1$,
\begin{align}
 \calR_{\rm cp}(n,m)
 &=\inf_\pi\left\{B_n^\pi+q_m(\nu_n^\pi)\right\},
 &B_n^\pi&=\E_\pi g(X_n),\label{eq:cp-exact}\\
 c_0m^{-2/d}&\le\calR_{\rm cp}(n,m)-V_0(u_0)
 \le C_0m^{-2/d},\label{eq:cp-main}\\
 c_0(M-n)^{-2/d}&\le\calR_{\rm aut}(n,M)-V_0(u_0)
 \le C_1(M-n)^{-2/d}.\label{eq:aut-main}
\end{align}
Here $q_m(\nu)=\inf_{z_1,\ldots,z_m}\int\min_j\norm{x-z_j}^2\nu(dx)$, and centroids may be chosen in $K$. Constants depend only on the fixed certificates, the domain and task, not on $n,m,M$. The autonomous class is empty when $M<n+1$.

The upper in~\eqref{eq:aut-main} is achieved to the stated bound by one Borel feedback controller with at most $M$ total labels. Its actions are Bellman-minimizing at its own actual retained barycenters. The lower is against every legal randomized Borel controller and is relative to the same unrestricted full-history optimum. Its programmed constants, workspace and finite-precision implementation are separately accounted for in Proposition~\ref{prop:numeric}.
\end{theorem}

\begin{remark}[The assumption hierarchy]\label{rem:hierarchy}
The proof has three levels. Raw affine response yields the exact calibration and Bellman-defect identity, Proposition~\ref{prop:identity}; no stability or geometry is used. Acquired density and terminal task curvature yield the all-policy lower. For the matching upper, acquired overlap and an interior buffer supply summable Bellman sensitivity, and the global grid makes the recursion executable. These are sufficient conditions, not an intrinsic necessary-and-sufficient classification. The overlap route does not require a same-report posterior contraction or a one-step Wasserstein contraction of the acquired kernel. Proposition~\ref{prop:expansive} gives a strict separation from the latter requirement.
\end{remark}

\begin{proposition}[A raw report-corridor certificate]\label{prop:corridor}
Suppose $C_0=D,C_1,\ldots,C_{b-1}$ are compact sets. For each $i<b$ and every action, there is a Borel report event $E_i^a$ with conditional probability at least $p_i>0$ at every incoming state in $C_{i-1}$, and its posterior lies in $C_i$. At the last step suppose
\[
 \mu_{s+b,a,u}\ge q\zeta_s
 \quad (u\in C_{b-1},\ a\in A)
\]
for $q>0$. Then~\eqref{eq:overlap} holds with
$\epsilon_*=q\prod_{i=1}^{b-1}p_i$, uniformly over feedback policies and their private randomizations.
\end{proposition}
\begin{proof}
At each conditioned history in the corridor, the action actually chosen has the stated lower event probability and maps into the next set. Successive conditional integration gives probability at least $\prod p_i$ of reaching $C_{b-1}$ along these events. For a nonnegative terminal test $f$, integration of the last-row lower on this subevent gives $\E f(X_{s+b})\ge q\prod p_i\int f\,d\zeta_s$. Randomized actions are mixtures of the same inequalities. The remaining paths retain their full nonnegative mass. Thus no conditional normalization or free rejection of failed corridor trials is used.
\end{proof}

For comparison, define $L_t=\Lip_K(V_t)$. Even without~\eqref{eq:overlap}, whenever these constants are finite the proof gives the weighted upper
\begin{equation}\label{eq:weighted}
 \calR_{\rm aut}(n,M)-V_0(u_0)
 \le C\inf_{\substack{m_t\ge1\ {
m integers}\\1+\sum_tm_t\le M}}
                 \sum_{t=1}^n L_t m_t^{-2/d}.
\end{equation}
No matching lower for an arbitrary nonsummable sequence $L_t$ is asserted. Theorem~\ref{thm:main} proves that actual overlap makes this right-hand side horizon-uniform, rather than assuming the desired continuation decay.

\section{Proof: acquired overlap, Bellman curvature and executable allocation}\label{sec:proof}
\subsection{Conditional calibration and exact defects}
The next proposition states the inherited conditional mechanism in its controlled form. It also records exactly where action optimality enters.
\begin{proposition}[Controlled calibration and Bellman defects]\label{prop:identity}
Let a legal finite controller have actual incoming barycenter $Q_{t-1}=\E[X_{t-1}\mid S_{t-1}]$. After its chosen action and report put
\[
 W_t=F_{t,A_t}(Y_t,Q_{t-1}),\qquad Q_t=\E[W_t\mid S_t].
\]
Then $Q_t=\E[X_t\mid S_t]$. For its best terminal readout, the exact risk decomposition is
\begin{equation}\label{eq:bellman-defect}
 R-V_0(u_0)=\sum_{t=1}^n\Delta_t+\sum_{t=1}^n J_t,
\end{equation}
where
\begin{align*}
 \Delta_t&=\E\left[\int V_t(w)\mu_{t,A_t,Q_{t-1}}(dw)
                         -V_{t-1}(Q_{t-1})\right]\ge0,\\
 J_t&=\E V_t(Q_t)-\E V_t(W_t)\ge0.
\end{align*}
An arbitrary terminal response adds its nonnegative excess over the conditional mean.

For completeness, set
$\Phi_{t,a}(u)=\int\norm{U_{t,a}(y,u)}^2/\ell_{t,a}(y,u)\,d\lambda_{t,a}$.
If $r_t=\E\norm{X_t-Q_t}^2$, the normalized three-defect identity is
\begin{equation}\label{eq:three-defect}
 r_t=\E[\Phi_{t,A_t}(X_{t-1})-\Phi_{t,A_t}(Q_{t-1})]
                 +D_{{\rm row},t}+D_{{\rm pool},t}.
\end{equation}
All terms are nonnegative. These identities require the raw affine response, not the hypotheses of Definition~\ref{def:class}.
\end{proposition}
\begin{proof}
Conditional on an incoming label and its newly chosen action, fresh action randomization is independent of the old physical state and history. Thus the conditional mean remains $Q_{t-1}$. Affinity in~\eqref{eq:raw} gives, by conditional integration,
\[
 \Pp(Y_t\in dy\mid S_{t-1},A_t)=\ell_{t,A_t}(y,Q_{t-1})\lambda_{t,A_t}(dy),
\]
\[
 \E[X_t\one_{\{Y_t\in dy\}}\mid S_{t-1},A_t]
     =U_{t,A_t}(y,Q_{t-1})\lambda_{t,A_t}(dy).
\]
Division off the null set proves $W_t=\E[X_t\mid S_{t-1},A_t,Y_t]$. The update coins are fresh, and all incoming edges into a new label are included; the tower property proves the assertion about $Q_t$. This proof uses measurable conditional integration rather than conditioning a continuous report on a positive-probability singleton.

Concavity of $V_t$ follows by taking the infimum of the affine risks of all fixed continuation trees with committed terminal responses. Jensen's inequality gives $J_t\ge0$; Bellman optimality gives $\Delta_t\ge0$. Adding the displayed definitions telescopes to $\E g(Q_n)-V_0(u_0)$, which is precisely the best label-based risk by~\eqref{eq:task}.

For~\eqref{eq:three-defect}, let $C_t=\E[W_t\mid S_{t-1},A_t,S_t]$. Two conditional Pythagorean identities give
\[
 \E\norm{W_t}^2-\E\norm{Q_t}^2
 =\underbrace{\E\norm{W_t-C_t}^2}_{D_{{\rm row},t}}
 +\underbrace{\E\norm{C_t-Q_t}^2}_{D_{{\rm pool},t}}.
\]
The raw identities identify $\E\norm{X_t}^2$ and $\E\norm{W_t}^2$ with the two perspective integrals in~\eqref{eq:three-defect}. Their difference is also
$\E\norm{X_t-\E[X_t\mid S_{t-1},A_t,Y_t]}^2$.
This proves its sign without differentiability. When differentiation is dominated, the unambiguous second-variation formula is
\begin{equation}\label{eq:hessian}
 \frac12D^2\Phi_{t,a}(u)[v,v]
 =\int\frac{\norm{A_{t,a}(y)v-(c_{t,a}(y)v)F_{t,a}(y,u)}^2}
                 {\ell_{t,a}(y,u)}\,d\lambda_{t,a}.
\end{equation}
It follows by differentiating the scalar quadratic perspective; otherwise its closed convex extension suffices. No covariance or Fisher metric has been assumed.
\end{proof}

\subsection{The new continuation mechanism}
\begin{lemma}[From acquired overlap to Bellman sensitivity]\label{lem:span}
Under~\eqref{eq:overlap}, with $\vartheta=1-\epsilon_*$,
\begin{equation}\label{eq:span}
 \osc_D V_t\le \vartheta^{\lfloor(n-t)/b\rfloor}\osc_D g,
 \qquad
 \Lip_K V_t\le r^{-1}\vartheta^{\lfloor(n-t)/b\rfloor}\osc_D g.
\end{equation}
For $\vartheta=0$, the expression is interpreted as one for exponent zero and zero otherwise.
\end{lemma}
\begin{proof}
For a bounded terminal continuation $f$, every $b$-step policy and initial state has terminal law $\epsilon_*\zeta_s+(1-\epsilon_*)\xi$ for a probability $\xi$. Its expectation lies in
\[
 \epsilon_*\int f\,d\zeta_s+
           (1-\epsilon_*)[\inf_D f,\sup_D f].
\]
Taking the infimum over policies preserves this same interval. Therefore the oscillation of the $b$-step Bellman transform is at most $(1-\epsilon_*)\osc_D f$. A shorter block cannot increase the range. Iteration proves the first inequality. This is the common-mass span contraction principle; its use here is on the acquired predictive kernels uniformly over feedback, not on a latent transition alone.

Let $v$ be concave on $D$, and let $x,y\in K$. Extend the line segment from $x$ backwards by length $r$, and from $y$ forwards by length $r$; both extension points belong to $D$. Monotonicity of the slopes of a concave function on this line bounds the slope on $[x,y]$ in absolute value by $\osc_Dv/r$. This proves the second inequality. No differentiability or choice of optimal action on a switching surface is required. In particular, one-step Wasserstein Lipschitz coefficients of the controlled acquired kernel never enter the argument.
\end{proof}

\begin{lemma}[Acquired density integrates nonsmooth curvature]\label{lem:curvature}
Let $v$ be concave and $L$-Lipschitz on a compact convex set $K\subset\R^d$ with nonempty interior, and let $\lambda$ be a probability on $K$ with density at most $\Lambda$. For each integer $m\ge1$ there is a global grid partition $\mathcal P_m$ with at most $m$ cells such that
\begin{equation}\label{eq:jensen-grid}
 \sum_{C\in\mathcal P_m}\lambda(C)v\left(\frac{\int_Cx\,d\lambda}{\lambda(C)}\right)
     -\int v\,d\lambda\le C_{K,d}\Lambda Lm^{-2/d}.
\end{equation}
Zero-mass cells are omitted. The grid does not depend on $v$ or $\lambda$; its conditional centroids do.
\end{lemma}
\begin{proof}
We supply the boundary-valid argument from the acquired-curvature method of \cite{GTFControl}. The convex function $f_0=-v$ has a convex $L$-Lipschitz extension to all $\R^d$: take the supremum of all supporting planes at interior points of $K$. All slopes have norm at most $L$; agreement on the interior and continuity give agreement on $K$. Let $\kappa=\Delta f$ be its positive distributional Laplacian measure.

For a grid cube $C$ of side $h$, centered at $z_C$, first mollify $f$ and subtract its tangent at $z_C$, obtaining a nonnegative convex function $w$ vanishing there. Put $R=2\sqrt d h$. The submean inequality, radial monotonicity and convex radial slopes give, for $x\in C$,
\[
 w(x)\le C_d R^{-d}\int_{B_R(z_C)}w
 \le C_d\operatorname{av}_{\partial B_R}w
 \le C_dR\operatorname{av}_{\partial B_R}\partial_rw
 =C_dR^{2-d}\int_{B_R(z_C)}\Delta f.
\]
In dimension one the boundary average is the average at the two endpoints. The tangent terms cancel at the conditional barycenter, and $w$ is nonnegative there. Passing through mollification and enlarging the ball slightly proves the cell bound
\[
 \int_C f\,d\lambda-\lambda(C)f(\E[X\mid C])
 \le C_dh^{2-d}\lambda(C)\kappa(\overline B_{3\sqrt d h}(z_C)).
\]
On a fixed bounded neighborhood of $K$, a smooth nonnegative cutoff $\chi$ gives
$\int\chi\,d\kappa=-\int\nabla\chi\cdot\nabla f\le L\int|\nabla\chi|$.
The enlarged grid balls have bounded overlap and $\lambda(C)\le\Lambda h^d$. Summation yields $C_{K,d}\Lambda Lh^2$. A bounding cube divided in each coordinate into $\lfloor m^{1/d}\rfloor$ intervals has at most $m$ cells and $h\le C_Km^{-1/d}$. This also covers $m=1$ after adjusting the fixed constant. The extension and its curvature measure are proof objects, not observer memory.
\end{proof}

\subsection{Implementation, lower bound and clock}
\begin{lemma}[Calibrated feedback implementation]\label{lem:implementation}
For any positive integers $m_1,\ldots,m_n$, there is a controller with disjoint phase alphabets of these sizes and one initial label, whose risk obeys
\begin{equation}\label{eq:implemented}
 R-V_0(u_0)\le C\Lambda\sum_{t=1}^n\Lip_K(V_t)m_t^{-2/d}.
\end{equation}
Its next-label rule is defined for all legal incoming labels and reports.
\end{lemma}
\begin{proof}
Inductively assign to every positive-mass label its actual conditional barycenter $q_s$. Choose a minimizing action in~\eqref{eq:bellman} at $q_s$; a finite action set and continuous Bellman transforms give a Borel selector, with a fixed tie rule. Compute $F_{t,a_s}(y,q_s)$ and send it to its cell in the single global partition $\mathcal P_{m_t}$. The next representative is the conditional mean of this unpooled value after integrating over \emph{all} incoming labels and reports entering that cell. It belongs to $K$ by convexity. Arbitrary legal representatives are assigned to unused labels.

Proposition~\ref{prop:identity} proves calibration for this very implemented law. Its row mixture has density at most $\Lambda$. Every $\Delta_t$ vanishes, and Lemma~\ref{lem:curvature} bounds $J_t$. Summation proves~\eqref{eq:implemented}. The only retained dynamic datum is the label; the calibrated centroid and selected action are programmed constants attached to it. There is no independent exact posterior state. The grid and update are global, so an unlikely but legal report does not escape the controller's definition.
\end{proof}

\begin{lemma}[All-policy acquired-mass lower]\label{lem:lower}
Every terminal acquired law of a full-history policy has density at most $\Lambda$. Encoding it into $m$ labels incurs squared distortion at least
\begin{equation}\label{eq:smallball}
 c_d\Lambda^{-2/d}m^{-2/d},\qquad
 c_d=\tfrac12(2v_d)^{-2/d},
\end{equation}
where $v_d$ is unit-ball volume. For every controller, the final term $J_n$ in~\eqref{eq:bellman-defect} satisfies the same bound with the number of terminal labels in place of $m$.
\end{lemma}
\begin{proof}
Condition before the last report on the actual incoming history and chosen action. The row bound in Definition~\ref{def:class} survives mixing, including private action randomness. For any $m$ centers, their radius-$r_0$ balls have total probability at most $\Lambda mv_dr_0^d$. Choose $r_0=(2\Lambda mv_d)^{-1/d}$; at least half the mass is outside. This proves~\eqref{eq:smallball}. Randomized assignment cannot improve the pointwise minimum over its finitely many readouts.

Since $g=c-\norm\cdot^2$ with affine $c$, $J_n=\E\norm{W_n-Q_n}^2$. Its actual unpooled law is also a mixture of rows at retained barycenters and chosen actions, hence has the same density bound. Repeated readouts only reduce the number of distinct centers. All previous defects and action slacks are nonnegative, which is why the converse compares with the full-history adaptive optimum rather than with a fixed controller.
\end{proof}

\begin{lemma}[Internal exact-deadline charge]\label{lem:clock}
Under the autonomous support condition, positive-probability label supports at cuts $0,1,\ldots,n$ of any legal internally clocked controller are pairwise disjoint. Consequently $M\ge n+1$ and the last cut has at most $M-n$ labels.
\end{lemma}
\begin{proof}
If one label occurred at cuts $s<t$, run a reference trace from that label using the controller's own action and update randomization and drawing each report from $\lambda_a$. Every finite actual future trace, conditional on any possible incoming physical history, is equivalent to this reference trace: iterated conditional densities are strictly positive and finite almost everywhere. The assertion is about positive-measure label edges, not individual report values. From the occurrence at $t$, halting after $n-t$ further calls has probability one. Equivalence forces the same event from the occurrence at $s$, contradicting the exact deadline. A halt label cannot be nonhalt at an earlier cut either. Each cut is nonempty. Conversely a deterministic chain with $n+1$ labels and arbitrary legal fixed actions realizes the deadline.
\end{proof}

\begin{lemma}[Integer allocation]\label{lem:allocation}
Put $p=2/d$, $a_t=r^{-1}\osc_Dg\,\vartheta^{\lfloor(n-t)/b\rfloor}$ and $A_n=\sum_ta_t^{1/(1+p)}$. For $k\ge1$ there are integers $m_t\ge1$ with
\[
 \sum_tm_t\le n+k-1,\qquad
 \sum_ta_tm_t^{-p}\le2^pA_n^{1+p}k^{-p}.
\]
Moreover $A_n\le b(r^{-1}\osc_Dg)^{1/(1+p)}/(1-\vartheta^{1/(1+p)})$, with the evident interpretation at $\vartheta=0$.
\end{lemma}
\begin{proof}
For $A_n>0$ set $w_t=a_t^{1/(1+p)}$ and
$m_t=1+\lfloor(k-1)w_t/A_n\rfloor$. When $k\ge2$, $m_t\ge(k-1)w_t/A_n$, and summation gives $A_n^{1+p}(k-1)^{-p}\le2^pA_n^{1+p}k^{-p}$. For $k=1$, use $\sum a_t\le A_n^{1+p}$. The geometric series proves the last assertion. If all weights vanish, one label per phase suffices. This counts a real phase label even at an insensitive cut; constants deteriorate explicitly as $\vartheta\uparrow1$.
\end{proof}

\begin{proof}[Proof of Theorem~\ref{thm:main}]
Conditional projection of the audit onto the full terminal history and then onto the checkpoint label proves~\eqref{eq:cp-exact}; measurable nearest-center quantization gives the infimum, and conditional centroids lie in $K$. The baseline of every policy is at least $V_0(u_0)$. Lemma~\ref{lem:lower} proves the checkpoint lower; a full-history Bellman-optimal policy followed by a global grid gives its upper. Weak continuity, finite actions and compact response sets give the needed Borel selectors by backward induction, or an arbitrarily small slack followed by a limit of the displayed risk bound.

For the autonomous upper put $k=M-n$. Lemmas~\ref{lem:span}, \ref{lem:implementation} and~\ref{lem:allocation} give the desired controller with $1+\sum m_t\le M$. The bound on $A_n$ is independent of $n$. For the lower, Lemma~\ref{lem:clock} and the final Jensen lower in Lemma~\ref{lem:lower} give $c_0(M-n)^{-2/d}$ for every controller. The clock lemma also proves the exact nonemptiness threshold. This establishes all statements without an interchange of minimax quantifiers or an oracle action sequence.
\end{proof}

In operational terms, the alphabet is allocated as follows:
\[
\begin{array}{c|ccc}
\text{cut}&0&1,\ldots,n-1&n\\\hline
\text{role}&\text{initial action}&\text{acquire, update, count}&\text{halt and respond}\\
\text{labels}&1&m_1,\ldots,m_{n-1}&m_n
\end{array}
\]
Thus $M-n$ is an upper bound on remaining terminal width after paying at least one label at each earlier cut. It is not an average width or a bit budget. A free time stamp or free terminal request changes this architecture and invalidates this clock subtraction.

\section{Active Gaussian observation, including expansive steps}\label{sec:classical}
Let the hidden state take values in $\{1,\ldots,q\}$, $q\ge2$. Fix an invertible column-stochastic matrix $P$ with $P_{ij}\ge\alpha>0$; necessarily $\alpha\le1/q$. At each step the controller chooses a sensor $a$ from a finite set. Conditional on the \emph{current} hidden state $i$, it receives
\[
 Y\sim N(m_i^a,\Sigma_a)\quad\text{in }\R^{q-1},
\]
then the hidden state makes the transition $P$. Each $\Sigma_a$ is positive definite and the vectors $m_i^a-m_q^a$, $i<q$, are linearly independent. No report is truncated. The protected terminal audit is the categorical vector of the post-transition hidden state; the loss is ordinary Brier squared error.

\begin{theorem}[Optimal finite-state active Gaussian sensing]\label{thm:gaussian}
Every such finite sensor family satisfies Theorem~\ref{thm:main} with $d=q-1$, $K=P\Delta_q$, and $D=\{u\in\Delta_q:u_i\ge\alpha/2\}$. In particular the matched laws hold against the full-history optimum over all adaptive sensor choices, for every positive invertible $P$, not only transitions satisfying a strong one-step contraction coefficient. Constants are uniform on compact parameter classes with fixed positive lower bounds on $\alpha$, singular values of $P$ and the likelihood-rank matrices, and covariance eigenvalues, and fixed upper bounds on means and covariances.
\end{theorem}
\begin{proof}
Equations in this section use Lebesgue report densities. Dividing all $f_i^a$ by one fixed everywhere-positive Gaussian reference density gives the probability-reference convention of Section~\ref{sec:experiments}, with the same normalized update and acquired law. Writing $f_i^a$ for the Gaussian densities, the raw instrument is
\begin{equation}\label{eq:gauss-raw}
 L_{a,y}u=P\diag(f_1^a(y),\ldots,f_q^a(y))u,\quad
 \ell_a(y,u)=\sum_i u_if_i^a(y),\quad
 F_a(y,u)=P\frac{(u_if_i^a(y))_i}{\ell_a(y,u)}.
\end{equation}
The terminal category tests separate beliefs; every finite executable test has an affine expectation in the initial belief, by conditioning on the hidden category and integrating its legal continuation. This proves the finite test closure directly. Mixture preparations realize $D$. Log odds give a continuous inverse report section on the interior of $K$ at each fixed phase; boundary mixtures are prepared directly. The predictive quotient therefore has the required Borel realization. Since $P\Delta_q\subset\{u_i\ge\alpha\}$, there is a positive relative Euclidean buffer between $K$ and the boundary of $D$.

Before applying $P$, put $w_i=u_if_i^a(y)/\ell_a(y,u)$ and $z_i=\log(w_i/w_q)$ for $i<q$. Then
\begin{equation}\label{eq:logodds}
 z_i=\log(u_i/u_q)+(m_i^a-m_q^a)^T\Sigma_a^{-1}y
 -\tfrac12\left((m_i^a)^T\Sigma_a^{-1}m_i^a
                      -(m_q^a)^T\Sigma_a^{-1}m_q^a\right).
\end{equation}
Under the actual report mixture, $z$ has a finite mixture of nonsingular Gaussian densities. The covariance determinant has a positive lower bound from the stated likelihood-rank determinant and covariance bounds. The means are uniformly bounded for $u\in D$. Consequently its density is at most $C\exp(-c\norm z^2)$, uniformly over $u,a$. The inverse logistic Jacobian in simplex coordinates is $(\prod_iw_i)^{-1}$; it grows at most $C\exp(C\norm z)$. Hence the acquired density of $w$ is bounded throughout the open simplex, including the tails approaching its boundary. The invertible affine map $P$ gives the required density bound on $K$.

On the fixed compact simplex section $W_*:=\{w_i\ge1/(2q)\}$, the same Gaussian formula has a strictly positive uniform lower bound, because its inverse log odds are bounded and all mixture parameters range in compact sets. Thus every acquired row at every $u\in D,a$ dominates a common positive multiple of uniform volume on $PW_*$. This is~\eqref{eq:overlap} with $b=1$. This lower is on \emph{actual report-induced density}; its probability is not normalized to one after selection.

The likelihoods are positive on all $\R^{q-1}$, with common null sets. Moreover
$\TV(\ell_a(\cdot,u),\ell_a(\cdot,v))\le\tfrac12\sum_i|u_i-v_i|$.
Pointwise continuity of~\eqref{eq:gauss-raw} and domination by $\sum_i f_i^a$ give weak continuity of the acquired kernels. Global grids and legal centroids follow from convexity of $K$. The categorical audit has $c=1$ and $g(u)=1-\norm u^2$. All assumptions of the same general theorem are verified.
\end{proof}

The acquired-overlap proof is different from the Hilbert-metric proof of passive posterior stability. For each fixed report, likelihood multiplication is a Hilbert isometry and a positive transition contracts. But feedback changes future report choices, and the relevant controlled Bellman comparison is not obtained by declaring the passive comparison optimal. The next proposition gives both an explicit failure of a sufficient one-step Wasserstein route and an actual action-switching effect.

\begin{proposition}[Expansive acquired transport and strict adaptive sensing]\label{prop:expansive}
Let $q=3$, $P=.9I+.1\one\one^T/3$, $\Sigma_a=I_2$. Sensor 1 has means $(10,0),(0,h),(0,-h)$ with $h=10^{-6}$, and sensor 2 interchanges hidden categories 1 and 2 in this list. This positive, full likelihood-rank model has the following properties.

For sensor 1 there are $u,v\in K$ such that
$W_1(\mu_{1,u},\mu_{1,v})>\norm{u-v}$, with Euclidean ground metric. Nevertheless Theorem~\ref{thm:gaussian} applies. At the last acquisition, the strictly optimal sensor is different on neighborhoods of $(1/2,1/4,1/4)$ and $(1/4,1/2,1/4)$. Both neighborhoods have positive acquired probability after any first sensor at an interior preparation.
\end{proposition}
\begin{proof}
We first record a perturbation estimate. For two joint laws of a finite hidden category and a common report, with joint total variation $\tau$, the laws of their posterior category vectors have $W_1$ distance at most $6\tau$. To see this, couple the common part of their report marginals. If their joint densities are $p_i(y),q_i(y)$, the common-density integral of the posterior $\ell^1$ difference is at most
$\int\sum_i|p_i-q_i|+\int|\sum_i(p_i-q_i)|\le4\tau$.
Unmatched report mass is at most $\tau$ and simplex diameter is $\sqrt2$. This proves the bound. Application of $P$ on the tangent space multiplies distances by $.9$.

At $h=0$, the report depends only on the partition $\{1\},\{2,3\}$. Its Bayes classification error is at most
$e=\Pp(N(0,1)>5)\le e^{-25/2}/(5\sqrt{2\pi})<3\cdot10^{-7}$.
Couple its posterior with the perfect-partition posterior. The expected difference is at most $2\sqrt2e$: if $p_y$ is the posterior probability of the first part, use $\E|p_y-\one_{\{i=1\}}|=2\E p_y(1-p_y)\le2e$.
Changing $h$ from zero changes the joint law by at most $h/\sqrt{2\pi}$. Thus, uniformly in the prior, the actual acquired row after $P$ is within
\begin{equation}\label{eq:gauss-perturb}
 2\sqrt2 e+6h/\sqrt{2\pi}<4\cdot10^{-6}
\end{equation}
of the ideal partition row in $W_1$.

Take $u=(1/3,1/3,1/3)$ and $v=u+.01(e_1-e_2)$. For the ideal row write
$w=(0,1/2,1/2)$ and
$z=(0,(1/3-.01)/(2/3-.01),(1/3)/(2/3-.01))$.
The common mass at $e_1$ is coupled identically. Optimal transport of the two remaining atoms costs
\[
 .9\left[.01\norm{e_1-w}+(2/3-.01)\norm{w-z}\right]
 =.009(\sqrt3+1)/\sqrt2>.01738,
\]
whereas $\norm{u-v}=.01\sqrt2<.01415$. Optimality follows from the two-by-two transport problem, or the triangle inequality after uncrossing any flow out of the common $e_1$ mass. The perturbation bound~\eqref{eq:gauss-perturb} leaves a strictly positive gap. Every actual likelihood-rank determinant is nonzero because $h>0$.

For the ideal sensors, the difference between expected squared posterior norms before $P$ is
\[
 \frac{2(p_1-p_2)p_3^2}{(1-p_1)(1-p_2)}.
\]
At $(1/2,1/4,1/4)$ it is $1/12$; interchanging the first two coordinates changes the sign. Applying $P$ multiplies this difference by $.9^2$, giving $.0675$. The function $\norm x^2$ is 2-Lipschitz on the simplex, so~\eqref{eq:gauss-perturb} preserves the strict comparison for the finite Gaussian model. Continuity gives neighborhoods. Formula~\eqref{eq:logodds} gives positive density on the entire interior of $K$, so both neighborhoods are actually encountered with positive probability. At horizons at least two a last-step feedback choice is therefore strictly better, conditionally on the preceding acquired state, than using either fixed sensor on both regions. This is a decision effect, not merely a parametrization of the same passive reports.
\end{proof}

\subsection{A fully specified binary benchmark}
The following passive subfamily is included to expose constants and a finite program without hiding a numerical optimization behind the theorem. Take two hidden categories, means $+1,-1$, variance one, and
$P=\left(\begin{smallmatrix}3/4&1/4\\1/4&3/4\end{smallmatrix}\right)$.
Write $q$ for the first-category probability. Then $D=[1/8,7/8]$, $K=[1/4,3/4]$, and
\begin{equation}\label{eq:benchmark-update}
 F(y,q)=\tfrac14+\frac{1/2}{1+((1-q)/q)e^{-2y}}.
\end{equation}
The exact categorical risk is $2\E(q_n-\widehat q_n)^2+B_n$. The logistic log odds have variance four and means bounded in absolute value by $2+\log7$. A direct maximization of the Gaussian density divided by the logistic Jacobian gives the conservative coordinate density bound
\begin{equation}\label{eq:benchmark-constants}
 \Lambda_0=\frac{28e^4}{\sqrt{2\pi}}<610,
 \qquad
 \epsilon_0=\frac1{\sqrt{2\pi}}
       \exp\left(-\frac{(2+\log21)^2}{8}\right)>0.
\end{equation}
Indeed $(1+e^z)^2/e^z\le4e^{|z|}$ and
$\sup_z\{|z|-(z-\mu)^2/8\}\le|\mu|+2$ prove the upper, including the $1/2$ transition scale. On $q'\in[3/8,5/8]$ the log odds have absolute value at most $\log3$; the inverse Jacobian is at least eight. Its length $1/4$ gives the displayed common acquired mass lower.

A scheduled three-report machine uses widths $10,20,30$, with one initial label, hence $M=61$ total states. At phase $t$ store only a midpoint from
$\{1/4+(j-1/2)/(2m_t):1\le j\le m_t\}$ and update by~\eqref{eq:benchmark-update} followed by nearest midpoint. Start with the known preparation $q_0$. For every report and every $q\in K$,
$|\partial_qF|\le2/3$; the grid error is at most $1/(4m_t)$. Consequently
\[
 |q_3-\widehat q_3|
 \le(2/3)^2/40+(2/3)/80+1/120=1/36,
 \qquad R-B_3\le1/648.
\]
The unrestricted 61-label checkpoint has upper excess $1/(8\cdot61^2)$; the acquired-density converse gives respectively
$1/(16\Lambda_0^2\cdot30^2)$ and $1/(16\Lambda_0^2\cdot61^2)$ as valid lower bounds. They are conservative constants, not equalities.

Every transition threshold is explicit: for a boundary $b\in(1/4,3/4)$, compare $y$ to
$\frac12\log\{[(b-1/4)/(3/4-b)](1-q)/q\}$.
There are $9+10\cdot19+20\cdot29=779$ such threshold constants and 30 rational terminal readouts. The experiment does not report a truncated cell. A finite-precision implementation must pay for representing these logarithms and handle a tail or boundary failure with Proposition~\ref{prop:numeric}. This benchmark is a legal construction with proved error bounds, not a numerical claim of attaining the exact finite-$M$ optimum.

\section{Noncommuting instruments with genuine feedback}\label{sec:quantum}
Let the physical Hilbert space have finite dimension $D_q\ge2$, and let $\mathcal H_0$ be its real space of traceless Hermitian matrices, with Hilbert--Schmidt norm. Put $d=D_q^2-1$. Let $\mu$ be uniform volume on $\{H\in\mathcal H_0:\norm H_{\HS}\le1\}$. Fix $0<\eta<1$, $0<\lambda<1$, and a finite family of unitaries $U_a$. At the start of each call the controller chooses $a$ from the past reports. The raw instrument is
\begin{equation}\label{eq:quantum-raw}
 \mathcal I_{a,H}(\rho)=T_a(K_H\rho K_H),\quad
 K_H=(I+\eta H)^{1/2},\quad
 T_a(\sigma)=\lambda U_a\sigma U_a^*
                   +(1-\lambda)\tr(\sigma)I/D_q.
\end{equation}
The weak measurement is followed by the chosen unitary and positive channel; it is not a free unitary chosen after seeing this call's report. Subsequent calls may react to that report. The report density is
$\ell(H,\rho)=1+\eta\tr(H\rho)$, which is genuinely state-dependent, and the effects generally do not commute.

The terminal task uses any fixed finite informationally complete POVM and its categorical Brier loss. Its probability vector is an injective affine chart of the density matrices. For instance, if $B_j$ is a Hilbert--Schmidt orthonormal traceless basis, then $\norm{B_j}_{\rm op}\le1$ and the $2d$ effects $(I\pm B_j)/(2d)$ are a POVM whose probabilities separate states. This is one randomized executable measurement, not a joint observation of incompatible observables. The classical controller has no uncharged quantum workspace.

\begin{lemma}[Positive-congruence inverse]\label{lem:congruence}
For every positive definite density matrix $\rho$, the map
\[
 E\longmapsto\frac{\sqrt E\rho\sqrt E}{\tr(E\rho)}
\]
after restricting to $E>0$, $\tr E=D_q$, is a smooth diffeomorphism onto the positive definite density matrices. Its inverse at $\sigma$ is
\begin{equation}\label{eq:inverse}
 B=\rho^{-1/2}(\rho^{1/2}\sigma\rho^{1/2})^{1/2}\rho^{-1/2},
 \qquad E=\frac{D_q B^2}{\tr B^2}.
\end{equation}
Its Jacobian and inverse Jacobian are bounded on compact positive subsets of these domains.
\end{lemma}
\begin{proof}
$B$ is the unique positive solution of $B\rho B=\sigma$, by taking the positive square root after conjugation by $\rho^{1/2}$. A normalized solution is therefore a positive scalar multiple of $B$, and the trace-$D_q$ constraint fixes the scalar in~\eqref{eq:inverse}. Matrix square root and inversion are smooth on the positive cone; the two displayed maps are smooth inverses. Compactness gives the derivative bounds and a positive minimum absolute determinant. Since $H$ is traceless, the effects $I+\eta H$ do lie in the required trace hyperplane. This is a finite-dimensional positive-matrix argument, not an unbounded-operator closure assertion.
\end{proof}

\begin{theorem}[Finite-state quantum feedback acquisition]\label{thm:quantum}
The instrument~\eqref{eq:quantum-raw}, with a fixed informationally complete terminal POVM, satisfies Theorem~\ref{thm:main} with $d=D_q^2-1$. It holds for every $0<\lambda,\eta<1$ and every fixed finite unitary action family. Thus the comparison is with the optimal full-history feedback choice of unitaries of the same experiment, not with a prescribed passive word. Constants may deteriorate near $\lambda=1$, $\lambda=0$, $\eta=0$, or loss of informational completeness.
\end{theorem}
\begin{proof}
Each raw map is completely positive and its integral is trace preserving: $\int H\,d\mu=0$ and $\tr T_a(\sigma)=\tr\sigma$. All finite future test probabilities are linear in the incoming density matrix; the terminal POVM separates it. Thus the predictive chart and unnormalized closure are obtained from the primitive instrument. Preparation of arbitrary density matrices in the sets below gives a Borel realization section. No spectral geometry is assumed as a resource law.

In density coordinates choose
\[
 D=\{\rho:\rho\ge(1-\lambda)I/(2D_q),\ \tr\rho=1\},\qquad
 K=\{\rho:\rho\ge(1-\lambda)I/D_q,\ \tr\rho=1\}.
\]
There is a positive buffer between $K$ and the relative boundary of $D$. Every update maps $D$ into $K$. The report density lies between $1-\eta$ and $1+\eta$, and
$\TV(\ell(\cdot,\rho),\ell(\cdot,\sigma))\le(\eta/2)\norm{\rho-\sigma}_{\HS}$.
The normalized map is jointly continuous. Lemma~\ref{lem:congruence}, the compact effect ball with $I+\eta H\ge(1-\eta)I$, and the invertible tangent map $\lambda U_a(\cdot)U_a^*$ imply a uniform bounded $d$-dimensional acquired density. Fixed affine POVM charts preserve these properties with their fixed norm and determinant constants.

It remains to derive feedback-uniform acquired overlap, rather than assume a transport coefficient. Write $\rho_* = I/D_q$. At $\rho_*$ the normalized weak measurement is exactly
$\rho_*+\eta H/D_q$. After $T_a$ it is $\rho_*+\lambda\eta U_aHU_a^*/D_q$. Hence its image contains the common ball
\[
 B_*:=B_{\HS}(\rho_*,\lambda\eta/(4D_q))
\]
with inverse effects of norm at most $1/4$. By the explicit inverse~\eqref{eq:inverse}, uniformly over this closed ball and the finite action family, there is $r_*>0$ such that every incoming $\rho$ with $\norm{\rho-\rho_*}_{\HS}\le r_*$ has an inverse effect in the ball of radius $1/2$ for every target in $B_*$. The Jacobian has a finite upper bound on this compact set and $\ell\ge1-\eta$. Thus every such acquired row dominates $q_*$ times uniform probability on $B_*$ for an explicitly defined positive $q_*$: take the minimum density given by~\eqref{eq:inverse} times $\operatorname{vol}(B_*)$, reduced to at most one if necessary.

Now restrict one report to $\norm H_{\HS}\le h$ with $\eta h\le1/2$. Functional calculus gives $\norm{K_H-I}_{\rm op}\le\eta h$. Since $\tr(K_H\rho K_H)\ge1/2$ and $\norm\rho_{\HS}\le1$,
\[
 \left\|\frac{K_H\rho K_H}{\tr(K_H\rho K_H)}-\rho\right\|_{\HS}
 \le8\eta h.
\]
For every chosen unitary the next distance to $\rho_*$ is therefore at most
$\lambda\norm{\rho-\rho_*}_{\HS}+8\lambda\eta h$.
Let $R=\sup_{\rho\in D}\norm{\rho-\rho_*}_{\HS}$. Choose $b\ge2$ and $h>0$ satisfying
\begin{equation}\label{eq:quantum-corridor}
 \lambda^{b-1}R\le r_*/2,\qquad
 \frac{8\lambda\eta h}{1-\lambda}\le r_*/2,\qquad
 h\le1,\quad\eta h\le1/2.
\end{equation}
Such choices exist for every fixed $\lambda<1$. Each quiet report has conditional probability at least $(1-\eta)h^d$. After $b-1$ such reports the state is in the $r_*$-ball, regardless of the adaptive actions. Proposition~\ref{prop:corridor} now gives
\[
 \epsilon_*=q_*\bigl((1-\eta)h^d\bigr)^{b-1}>0.
\]
Every quiet report and every unsuccessful corridor attempt remains physically counted. This derives the required overlap from actual kernels even when the improving channel is arbitrarily weak. Common report support and all remaining hypotheses were already verified. The theorem follows.
\end{proof}

\begin{proposition}[A strictly feedback-dependent qubit task]\label{prop:qubit-switch}
For $D_q=2$, the preceding theorem contains a fixed physical detector for which the optimal last unitary switches on a positive-mass acquired surface. Let $n_1,\ldots,n_4$ be tetrahedral unit Bloch vectors, let
\[
 a_1=\tfrac12,\quad a_2=a_3=a_4=\tfrac16,\quad t=\tfrac1{16},
 \qquad E_i=a_iI+t\,n_i\cdot\sigma.
\]
Use the terminal POVM $(E_i)$ and actions given by the identity and a $\pi$ rotation about an axis orthogonal to $n_1$. The two actions are strictly ordered with opposite signs according to $n_1\cdot r$, where $r$ is the current Bloch vector.
\end{proposition}
\begin{proof}
Each $E_i$ is strictly positive because $a_i>t$, and they sum to $I$. Since $\sum n_i=0$ and $\sum n_in_i^T=(4/3)I$, their outcome probabilities $p_i=a_i+t n_i\cdot r$ separate all three Bloch coordinates. Their Brier Bayes risk is
\begin{equation}\label{eq:qubit-risk}
 g(r)=\tfrac23-\tfrac{2t}{3}n_1\cdot r-\tfrac{4t^2}{3}\norm r^2.
\end{equation}
The weak-effect report ensemble is rotationally invariant. For a fixed eigenvalue radius $z\in[0,\eta/\sqrt2]$, integration over directions of the unnormalized L\"uders map gives mean Bloch multiplier $(1+2\sqrt{1-z^2})/3$. Radial averaging gives a constant $\gamma>1/3$ with $\E_\ell[r_{\rm post}]=\gamma r$. This follows directly by resolving a vector into the effect direction, whose multiplier is one, and its two transverse directions, whose multiplier is $\sqrt{1-z^2}$.

After the chosen unitary and depolarization, the linear term in~\eqref{eq:qubit-risk} is $-(2t/3)\lambda\gamma n_1\cdot R_ar$. The expected squared-norm term does not depend on the post-measurement rotation. The difference between identity and the flipping action is exactly
$-(4t/3)\lambda\gamma n_1\cdot r$.
The acquired row from the maximally mixed preparation has positive density on a ball about the origin, so both signs occur with positive probability. The control aligns a partially observed state with a fixed asymmetric detector; it is not a change of notation for an invariant task. The detector remains protected until commitment, and no classical controller is granted additional quantum memory.
\end{proof}

\section{Visible changing-rank acquisition}\label{sec:singular}
The theorem also admits a singular acquired support without replacing real mixture weights by conditional probabilities of one. We give both an abstract finite-stratum statement and a raw controlled construction.

\begin{theorem}[Actual-mass stratified transfer]\label{thm:strata}
Replace the density part of Definition~\ref{def:class} by the following finite visible decomposition. Every acquired row has a current report mark $j\in\{1,\ldots,J\}$ with fixed probabilities $w_j>0$, $\sum_jw_j=1$. Conditional on that mark its predictive value lies in a compact convex set $K_j\subset\operatorname{int}D$ of affine dimension $d_j\ge1$, with nonempty relative interior and relative density at most $\Lambda_j$. All $K_j$ have a common positive buffer in $D$. Their convex hull lies inside $D$. The actual kernels, including their marks, satisfy feedback-uniform overlap and common actionwise support as before. Define
\begin{equation}\label{eq:xi}
 \Xi_w(m)=\sum_{j=1}^Jw_jm^{-2/d_j}.
\end{equation}
Then, for the same terminal quadratic task and all admissible budgets,
\[
 \calR_{\rm cp}(n,m)-V_0(u_0)\asymp\Xi_w(m),\qquad
 \calR_{\rm aut}(n,M)-V_0(u_0)\asymp\Xi_w(M-n).
\]
Constants may depend on the fixed component weights, dimensions, geometry and overlap. No optimal allocation over varying weights is claimed.
\end{theorem}
\begin{proof}
The current mark has the displayed probability at every incoming retained state and action. Hence conditioning on the last mark and applying the relative-volume union-of-balls argument gives, for every global codebook of $m$ centers, at least
$\sum_j c_jw_jm^{-2/d_j}$ distortion. A center outside the affine span of $K_j$ can only decrease its distance to that stratum when orthogonally projected; therefore the lower also holds for unrestricted global centers. All-policy and clock arguments are unchanged.

The overlap proof and concavity give a common bound
$\Lip_{K_j}V_t\le C\vartheta^{\lfloor(n-t)/b\rfloor}$ on every stratum. For the upper, keep the current mark in the outgoing cell label and partition each $K_j$ in its own affine coordinates. Conditional calibration then gives representatives in that same convex $K_j$, even after pooling all predecessor rows. Lemma~\ref{lem:curvature} in dimension $d_j$ gives the weighted row loss $C\vartheta^{\lfloor(n-t)/b\rfloor}\sum_jw_jm_{t,j}^{-2/d_j}$. If a phase width is at least $2J$, equal allocation with $m_{t,j}=\lfloor m_t/J\rfloor$ has this order. If it is smaller, one common label may merge marks; its Jensen loss is at most $\osc_DV_t$, which is bounded by a fixed multiple of the same profile for $m_t<2J$. Thus an extra $J$-state minimum is not imposed at every phase.

Use a geometric integer allocation $m_t=1+\lfloor(k-1)\widetilde w_t/\sum\widetilde w_s\rfloor$ with $k=M-n$ and $\widetilde w_t=\vartheta^{\lfloor(n-t)/b\rfloor/(1+p_*)}$, $p_*:=\max_j2/d_j$. For each $p_j=2/d_j\le p_*$, summing the geometric tail gives $\sum_t\vartheta^{\lfloor(n-t)/b\rfloor}m_t^{-p_j}\le C_jk^{-p_j}$, including $k=1$. This proves the autonomous upper. A terminal stratified code after a full-history optimal policy gives the checkpoint upper. The unconditional weights $w_j$ remain present throughout.
\end{proof}

\begin{corollary}[Raw rank-changing active sensors]\label{cor:rank-raw}
In the Gaussian controlled experiment, let the report first choose a visible mark $j$ with probability $w_j$, independently of the hidden category, then produce the chosen Gaussian sensor report, and then apply
\[
 P_j=(1-\lambda)\one\one^T/q+\lambda R_j,
 \qquad 0<\lambda<1,
\]
where $R_j$ is column stochastic of rank $d_j+1$. Suppose one component is $R_1=I$, with $w_1>0$, and each other $d_j\ge1$. The raw instrument is
$w_jP_j\diag(f_i^a(y))\,dy$ with report $(j,y)$. It satisfies Theorem~\ref{thm:strata} with $K_j=P_j\Delta_q$ and its stated actual dimensions. For $q=3$, taking the columns of one $R_j$ to be $e_1,e_1,e_2$ produces a one-dimensional segment in addition to the two-dimensional component.
\end{corollary}
\begin{proof}
The raw map is positive and preserves total mass after summing marks. All $P_j$ have entries at least $(1-\lambda)/q$, so the outer preparation domain with half that floor provides the buffer. The Gaussian pre-transition posterior $w$ has uniformly bounded full-dimensional density by~\eqref{eq:logodds}. Pushing it through the rank-$d_j$ tangent map $P_j$ gives bounded relative density: choose orthonormal coordinates splitting its kernel and complement, integrate over the bounded simplex fibers, and divide by the product of its nonzero singular values. These are fixed finite positive constants. This also proves the declared affine dimensions rather than reading them from ambient coordinates.

The full-rank branch has the common acquired density lower on a fixed compact simplex section from Theorem~\ref{thm:gaussian}. Multiplying by its true probability $w_1$ gives one-step feedback-uniform overlap. Common support holds on the whole marked report space because every $w_j$ and Gaussian density is positive. The other primitive verifications remain valid without requiring invertibility of the lower-rank transitions. All branches are paid physical outcomes. The mark weights refer to the \emph{current} report branch, not to persistence of a rank chosen at an earlier phase.
\end{proof}

For a purely atomic component, a finite number of reserved centers gives zero component distortion once its support is covered; its correct term is its finite-support quantization profile, not $m^{-2/0}$. We do not use this convention to claim a uniform theorem over infinitely many strata or vanishing component masses. The full-rank/pure-state quantum theorem in \cite{GTFBayes} is preserved separately with its original curvature assumptions; the controlled affine-stratum result above neither deletes nor silently generalizes that different orbit geometry.

\section{Causal composition, precision and an attained joint law}\label{sec:resources}
\subsection{Typed causal simulation}
A simulator from a source experiment to a target experiment is parameter independent and has a report transformation, a strategy lift specifying legal source actions for each target action, a charged internal state, a clock/call map, and a terminal task readout. Action compatibility includes the causal order of interventions. Its discrepancy is measured on the complete stopped and scored law, including the target controller's fresh randomization, committed response and protected audit, not only on report marginals.

\begin{proposition}[Same-task simulation composition]\label{prop:simulation}
Suppose a simulator uses at most $R$ internal labels, maps at most $n$ target calls to at most $c(n)$ source calls, and couples the complete stopped/scored laws with failure probability at most $\varepsilon$. Suppose the common loss takes values in $[0,L_{\max}]$. Every target controller with $M$ labels then induces a source controller with at most $RM$ labels and risk at most its target risk plus $L_{\max}\varepsilon+e_{\rm out}$, where $e_{\rm out}$ bounds any task-readout mismatch on the coupled event. Simulator program, workspace, physical time and calibration are separately charged. Two compatible simulators compose with product state counts, composed call maps, and error at most the sum of their complete-law coupling errors.
\end{proposition}
\begin{proof}
The source retains the ordered pair of simulator and target labels. The strategy lift reads this pair and produces only the currently permitted source action sequence; its internal clock is among the $R$ states. The simulator supplies the target report and the target update uses its own coupled fresh coin. Induction on completed calls couples the stopped target execution and protected audit. Loss agrees up to $e_{\rm out}$ on success and differs by at most $L_{\max}$ on failure. The union bound gives the composed discrepancy, and the paired state and composed physical call counts are literal. A one-way simulator gives this one-way implementable upper only; a reverse minimax inequality requires a reverse comparison or an independent source lower bound. No TV-to-large-deviation or finite-time-to-infinite-horizon inference is made.
\end{proof}

\subsection{Finite-bit execution is an additional assertion}
\begin{proposition}[Certified finite precision]\label{prop:numeric}
Consider an ideal calibrated controller in Lemma~\ref{lem:implementation}, possibly with action slacks summing in expectation to at most $\gamma$. Suppose it has a finite certificate for all incoming conditional-centroid integrals and action values, primitive evaluation, and a finite arithmetic implementation of its global cell rules. At phase $t$, assume primitive coupling failure probability $\xi_t$, arithmetic or tail failure probability $\zeta_t$, and a bound $h_t$ on the pre-quantization error outside these failures. Suppose the ideal actual unpooled row mixture obeys the \emph{unconditional} boundary bound
\[
 \Pp\{W_t\text{ is within }h_t\text{ of a cell boundary}\}\le\omega_t(h_t).
\]
If every rounded terminal response is within $\delta_{\rm out}$ of its programmed centroid, then a finite-bit implementation with the same dynamic labels and fixed action table has
\begin{equation}\label{eq:numerics}
 R_{\rm bit}\le V_0(u_0)+C(M-n)^{-2/d}+\gamma+\delta_{\rm out}^2
       +L_{\max}\sum_{t=1}^n[\xi_t+\zeta_t+\omega_t(h_t)].
\end{equation}
For the full-dimensional grids, $\omega_t(h)\le C\Lambda h m_t^{1/d}$ while $h$ is below the cell scale. Analogous relative boundary certificates apply to separately labeled strata.
\end{proposition}
\begin{proof}
Use the same ideal action table in both programs; a Bellman tie is handled offline by a fixed action with certified slack, not by assuming that close beliefs choose the same action. Define an intermediate reference execution using the ideal labels but the rounded terminal response table. Conditional calibration makes its risk exceed the ideal risk by exactly the mean squared readout rounding error, at most $\delta_{\rm out}^2$. This comparison is unconditional, so no invalid centering on the event of previous numerical success is required.

Couple the physical primitives and private coins of this reference execution and the finite implementation until the first label mismatch. Their fixed action rows then agree. A first mismatch can occur only on a primitive or arithmetic failure or when the ideal value lies in the boundary strip. For each phase this first-mismatch event is contained in the corresponding event on the \emph{unconditional ideal execution}; summing these probabilities is valid without conditioning a row law on past successful arithmetic. Until mismatch the rounded readout tables also agree. The complete loss bound and the union bound prove~\eqref{eq:numerics}. The action-slack term follows from~\eqref{eq:bellman-defect} for the actual ideal action table and its own forward-calibrated centroids. The grid boundary estimate follows by integrating the density over strips of total volume at most $C h m_t^{1/d}$.
\end{proof}

For $p$-bit constants a direct representation stores $O(dpM+M\log|A|+M\log M)$ bits for centroids, actions and state routing, in addition to the read-only grid and primitive programs. It uses the evaluator's declared report buffer and arithmetic workspace; none is identified with $\log_2M$. Offline Bellman evaluation and forward integration can be expensive and are not counted as zero time. The proposition assumes their certificates, rather than deducing computability from abstract Borel existence. For Gaussian evaluation a charged tail event may bound finite arithmetic; the physical experiment still reports the untruncated real value. A discrete rounded report is not close in TV to that continuous report merely because its numerical error is small.

\subsection{A protected post-control bit}
A bit acquired \emph{before} a controlled experiment can be encoded into the physical state through later actions. The retained-bit lower for a prescribed passive core in \cite{GTFBayes} cannot therefore be transplanted to controlled sensing without a new argument. The following ordering closes that side channel while keeping all resources and the same task explicit.

After preparation and the $n$ controlled core reports, prepare an independent fair $U\in\{-1,1\}$ and make one paid report $B$: it equals $U$ with probability $1-a$ and is an erasure symbol with probability $a$. No further core action is allowed. The observer then commits its response in $D\times[-1,1]\times[-\delta,\delta]$. The protected terminal audit scores the core target, $U$, and an unobservable fixed calibration offset $\theta\in\{-\delta,\delta\}$ by the sum of their squared losses. The program is independent of $\theta$. Write the corresponding minimax risks as $\calR^{\rm joint}$, and compare the bit-report interface with an exact-bit target that has the identical core, action order and audit.

\begin{theorem}[Attained acquisition--memory--calibration--deficiency law]\label{thm:joint}
The source-to-exact-bit causal deficiency in this protected interface is exactly $\varepsilon=a/2$. With $N=n+3$ physical calls and $n\ge1$, the checkpoint satisfies
\[
 \calR^{\rm joint}_{\rm cp}(N,m,\delta,\varepsilon)-V_0(u_0)
       \asymp m^{-2/d}+\delta^2+\varepsilon.
\]
The autonomous class is nonempty exactly when $M\ge n+2=N-1$, and on that class
\begin{equation}\label{eq:joint}
 \calR^{\rm joint}_{\rm aut}(N,M,\delta,\varepsilon)-V_0(u_0)
       \asymp (M-N+2)^{-2/d}+\delta^2+\varepsilon.
\end{equation}
Constants depend only on the fixed core certificates and task normalization, not on $n,M,a,\delta$ in a bounded declared calibration range. Replacing the core by Theorem~\ref{thm:strata} replaces its power by $\Xi_w$. These are attained errors of this experiment, not lower bounds inferred from arbitrary allowed simulator or numerical tolerances.
\end{theorem}
\begin{proof}
On erasure a stateless simulator guesses a fresh fair sign; otherwise it reports the observed sign. Its disagreement with the protected true $U$ is exactly $a/2$, and it reproduces all other parts of the target. Conversely the target has zero probability that its early completed bit report differs from the protected terminal $U$. In the source, conditional on an erasure and the entire core history, $U$ remains fair and inaccessible; no causal simulator can reduce disagreement below $a/2$. It cannot call the audit before completing the bit report or make an uncharged repeated bit acquisition. This proves exact deficiency.

The minimum bit squared risk from the full available information is $a$, attained by the responses $1,-1,0$ to the three symbols. The unobservable calibration contributes at least $\delta^2$ under the common equal prior on its two signs and is attained by response zero. These priors are independent of the core. The core lower from Theorem~\ref{thm:main} thus adds to these two lower terms, rather than requiring incompatible least-favorable alternatives.

For an autonomous observer, let $m_n$ be its number of labels immediately after the core and before the bit call. The preceding $n$ phase supports are disjoint by the positive-support trace argument. A bit-action label cannot serve at a core phase because the action type and prescribed order are fixed by the label; a halt label cannot serve before the bit. Hence
\[
 m_n\le M-n-1=M-N+2.
\]
The bit is prepared after every core action, so conditional on this core label its report and coins are independent of the core state. Processing them cannot reduce the core risk below the conditional-mean risk at that cut. The all-row final-core Jensen lower gives $c m_n^{-2/d}$. The complete bit and calibration lower bounds remain valid simultaneously. This proves~\eqref{eq:joint}'s lower and the analogous checkpoint lower. A chain of $n+2$ labels proves the exact feasibility threshold; its readout may be coarse.

For the upper, use the core calibrated controller, retain its terminal label during the bit call, and then create three terminal copies for the bit responses $1,-1,0$. If the core phase widths are $m_t$, the total is
\[
 1+\sum_{t<n}m_t+4m_n.
\]
For large enough $k=M-n-1$, the integer allocation proof applies with phase costs $c_t=1$ for $t<n$ and $c_n=4$, and with available extra budget $k-3$. The geometric sum remains bounded, since only one cost changed by a factor four. For $k\ge6$, $k-3\ge k/2$, giving core excess $Ck^{-2/d}$. For the finitely small values $1\le k<6$, a one-label-per-required-phase chain has bounded task loss, absorbed by a fixed multiple of $k^{-2/d}$. Use calibration response zero. The bit loss is $a=2\varepsilon$. A checkpoint uses three copies of a core code and the same small-budget convention. This proves both upper bounds. In the stratified case the same allocation is applied to the fixed finite collection of powers.
\end{proof}

The location of the bit call is part of this theorem. It neither overwrites the earlier passive retained-bit theorem nor silently changes its two-lane clock charge. In particular the physical count $N$ includes preparation, all $n$ controlled acquisitions, the separate bit report, and the audit; the audit is not an observation the controller may use to choose its committed response.

\section{Relation to neighboring theories and retained results}\label{sec:comparison}
Theorems~\ref{thm:main}, \ref{thm:gaussian} and~\ref{thm:quantum} concern an optimized feedback experiment, not merely an observer of a prescribed passive word. Their memory object is a total finite alphabet including actions, phases and stopping. This distinction matters both from the earlier passive theorem and from the following established lines of work.

Value-directed belief approximation and compression were developed by Poupart and Boutilier \cite{PB2000,PB2001,PB2002}. In particular, the lossless linear theory in \cite{PB2002} identifies invariant subspaces containing the reward vector, so that a compressed state predicts its own future and policy value; it also gives lossy value-error bounds. A finite executable affine closure is conceptually adjacent to that sufficiency principle, and is not advertised here as a newly discovered idea. The additional question addressed by Theorem~\ref{thm:main} is the error of an \emph{actually implemented hard-label recursion}, with all phases charged, and the converse over every randomized Borel controller of the same task. A finite-dimensional real linear compression is not itself a finite total label budget.

Bayes-risk error and Bregman divergence are linked explicitly by Varshney \cite{Varshney2011}; proper scoring geometry and convex perspectives provide broader context \cite{Gneiting2007,Boyd2004}. Our squared-loss Jensen identities are instances of this geometry. What must additionally be proved is causal calibration across all incoming rows, action slack relative to the true Bellman optimum, and an attained all-state lower. None follows merely from identifying a static Bregman divergence. The normalized three-defect identity of \cite{GTFBayes}, restated in Proposition~\ref{prop:identity}, is an inherited dynamic contribution, not a new static quantization principle.

Finite-state vector coding \cite{Foster1985}, optimal sequential vector quantization \cite{BMT2001}, and recursive marginal quantization \cite{Pages2015} are direct predecessors of finite-state and recursive approximation. Information-rate formulations such as sequential Gaussian rate--distortion \cite{Tanaka2017} charge a different resource. A source encoder observing the current source, a recursive codebook approximation, and a controller observing only noisy reports have different information patterns. The explicit $M-n$ constraint is the effect of the internal exact deadline in this architecture, not a replacement for their coding theorems.

The common-mass contraction in Lemma~\ref{lem:span} belongs to the classical ergodic-coefficient tradition \cite{Dobrushin1956}. Nonlinear-filter stability via projective metrics is also mature \cite{Atar1997}. Kara and Y\"uksel \cite{KY2022} prove near-optimal finite-memory feedback policies using filter stability, while Demirci, Kara and Y\"uksel \cite{DKY2024} establish average-cost optimality and approximation consequences from contraction-based regularity. Those results include control and infinite-horizon assertions that are not consequences of our finite terminal theorem. Here an actual acquired-overlap certificate, uniform over controlled report corridors, replaces the one-step Wasserstein sensitivity hypothesis used in the retained continuation-curvature argument \cite{GTFControl}. Proposition~\ref{prop:expansive} proves failure of that particular one-step sufficient hypothesis; it does not claim failure of every stability theorem or every possible multi-step metric.

Kim's finite-memory belief approximation \cite{Kim2026} measures truncated input--output histories and develops a policy-conditional Wasserstein comparison. Its central cost comparison evaluates a true-belief and a substituted-belief cost along one closed-loop execution. Our risk comparison instead permits the finite controller to change subsequent actions and compares its true scored execution with the unrestricted full-history optimum. D\'iaz-Infante Velasco, Garcia and Minj\'arez-Sosa \cite{DGM2026} formulate normalized class-restricted filters and discounted control for finite partition observations, with class-dependent approximations retaining atomic masses. Their observation-class complexity is not the total number of a controller's phase and halt states. The present actual-stratum theorem charges a visible mark when retained and proves a last-row converse; it does not identify arbitrary partition filters with our finite affine class. These recent preprints are treated as such, not as established priority certificates.

\subsection{What is unified and what remains independent}
Conditional pooling, Bellman defects, the boundary-valid curvature estimate, the phase allocation and the original positive-congruence calculation are retained results \cite{GTFControl,GTFAllocation,GTFBayes}, explicitly reproved where needed. The power $2/d$, the clock subtraction, and the elementary hidden-sign and erasure constructions are not first claimed in this article. The present theorem supplies a new sufficient continuation mechanism derived from feedback-uniform acquired mass, verifies it for active Gaussian sensing and feedback noncommuting instruments, and obtains their full-history-optimal hard-state curves. The rank-changing construction and the post-control bit theorem preserve the actual mass and the causal order needed by those claims.

Earlier exact noncontracting orbit and temporal laws, unknown-fixed-model common-program realization results, and the independent ordered-measurement, finite-action and streaming-space theories keep their own hypotheses. They are neither deleted nor treated as unproved premises of the present theorem. The complete reviewed Bayesian article is retained as a separate companion; an editor need not read a stack of historical manuscripts to verify the native proofs above.

The unresolved mathematical boundary is specific. Feedback-uniform acquired overlap is sufficient, not near-necessary. We do not prove matching converses for every nonsummable controlled continuation, unknown-kernel learning or regret, arbitrary nondominated supports, an infinite test carrier, or all degenerating and countably stratified laws. Exact deadlines are not optimized stopping times. Program description, workspace, planning time and minimal simulator state do not yet have a complete matching joint frontier. The current contribution advances the original experiment--geometry--resource problem on a broader controlled class; it does not turn that mother problem into a universal classification by terminology.

\appendix
\section{Compatibility, stability contrasts and interface checks}\label{sec:compatibility}
\subsection{Why local static equality is not global causal attainment}
If an incoming row has law $\mu_s$, any $m$-label randomized assignment has within-row distortion at least $q_m(\mu_s)$. In~\eqref{eq:three-defect}, equality in this row bound and vanishing pooling require, respectively, rowwise optimal quantization and agreement of all incoming centroids assigned to a common outgoing label. These are \emph{local} conditions. They do not assert the existence of one recursion satisfying them through every preceding information cut.

\begin{example}[Two explicit obstructions]\label{ex:compatibility}
Two equally likely rows $\delta_0$ and $\delta_1$ each have $q_1=0$. Merging them into one outgoing label has pooling distortion $1/4$, despite rowwise optimality. For an obstruction at an earlier cut, acquire two independent fair signs $B_1,B_2$, and only afterwards receive an independent fair index $I\in\{1,2\}$ requesting $B_I$. A two-label checkpoint after $I$ can answer exactly. A two-label state before $I$ cannot: its two possible readout pairs are a two-center quantizer for the four equally likely vertices of $\{-1,1\}^2$. The optimal total squared distortion is one, by checking adjacent two--two partitions versus one--three partitions; division by the independent index gives risk at least $1/2$. A randomized encoder cannot improve the pointwise minimum over its two response pairs. Thus a small final answer alphabet is not a global recursive memory certificate.
\end{example}

A full-dimensional terminal law is also insufficient for a uniform row condition. In the neutral circle-rotation experiment, a full-history predictive vector can be uniform on the unit circle while a one-label retained barycenter is zero. Its subsequent rotation row is then $\delta_0$. This is why the lower theorem quantifies over all legal retained barycenters, not merely a typical full-history law. The example is outside Definition~\ref{def:class}, and is used as an assumption witness, not a realization of its density theorem.

\subsection{Integrated curvature and same-report stability are different}
For comparison with the retained Bayesian theorem, call (C) the existence of block-summable upper curvature coefficients for the perspective potential, and call (P) block-summable pathwise Lipschitz coefficients for every same-report normalized update in a legal metric. Neither is silently substituted for the new acquired-overlap condition.

\begin{example}[(C) fails while (P) holds]\label{ex:Cfails}
A binary sign has prior mean $u$, a noisy sign report has density $(1+\eta yu)/2$, and a positive symmetric hidden transition multiplies the posterior mean by $\lambda$. Thus
$F(y,u)=\lambda(u+\eta y)/(1+\eta yu)$, $y\in\{-1,1\}$. Direct differentiation gives
\[
 \tfrac12\Phi''(u)=
 \lambda^2(1-\eta^2)^2\frac{1+3\eta^2u^2}{(1-\eta^2u^2)^3}.
\]
At $\lambda=\eta=u=.99$ this is greater than $24$. The homogeneous uniform (C) coefficient therefore cannot have summable products. In the metric $2|\operatorname{arctanh}u-\operatorname{arctanh}v|$, the report adds a constant to the arctanh coordinate, while $u\mapsto\lambda u$ has derivative ratio
$\lambda(1-u^2)/(1-\lambda^2u^2)\le\lambda<1$.
Thus (P) holds on the positive post-transition interval. The example separates stability routes; its finite report alphabet does not satisfy the full-dimensional row-density hypothesis.
\end{example}

\begin{example}[(P) fails while (C) holds]\label{ex:Pfails}
With probability $\omega\in(0,1)$ a marked quantum branch does nothing and reports its mark. Otherwise it performs the covariant rank-one instrument
$D_qP_\psi\rho P_\psi\,d\nu(\psi)$, where $\nu$ is uniform projective measure. Use the centered Hilbert--Schmidt predictive chart. The pure posterior has constant squared norm $1-1/D_q$, so the integrated potential is
\[
 \Phi(\rho-I/D_q)=\omega\norm{\rho-I/D_q}_{\HS}^2
                         +(1-\omega)(1-1/D_q).
\]
Its curvature coefficient is exactly $\omega<1$. On the identity branch every nondegenerate metric has same-report Lipschitz coefficient at least one, at every step. Hence no uniform pathwise (P) product contracts. All branch masses are real and positive; the closed perspective covers possible zero likelihood at an individual pure state. This finite-rank example is a contrast for the two sufficient routes, not an assertion of the interior/density hypotheses of Theorem~\ref{thm:main}.
\end{example}

\subsection{Clarifications for the retained passive companion}
The passive Gaussian convention in the retained Bayesian article is report first, positive transition second; its constants include the determinant lower of the affine likelihood-rank matrix and $\alpha\le1/q$. Its positive-cone contraction constants depend on $\lambda,D_q$, and its bounded-gap extension starts in a declared interior domain before the first improving step. In a projective-orbit small-ball estimate, a ball centered off the orbit either misses it or is contained, on that orbit, in a ball of twice the radius centered at an orbit point. Finite-stratum equal allocation is a sufficient fixed-weight upper, not a weight-optimal allocation theorem.

In that companion's \emph{passive} early-bit extension, the core remains independent of the bit even after conditioning on all its allowed reports and its fixed action word. A completely forgotten fair sign has squared loss one for responses in $[-1,1]$. The shrinkage response $(1-a)b$ in its two-lane construction has risk $2a-a^2$; with $\varepsilon=a/2$ this is at most $4\varepsilon$. If $Q=M-2n-1\ge2$, its core budget $M_0=1+\lfloor(M-3)/2\rfloor$ obeys
$M_0-n=\lfloor Q/2\rfloor\ge Q/3$; the remaining bounded values use its stated coarse construction. These are clarifications of the retained interface. They do not transfer the pre-core bit lower to arbitrary feedback actions. Theorem~\ref{thm:joint} instead proves the controlled post-core version with a different, explicit state count.

All finite-time simulation comparisons include the committed answer, target coins and protected audit. Loss bounds retain their physical-dimension normalization. These conventions rule out a free stored response, a latent external clock, or a simulator that reads the audit before generating its preceding report.

\clearpage\begin{thebibliography}{99}
\bibitem{Atar1997} R. Atar and O. Zeitouni, \emph{Exponential stability for nonlinear filtering}, Ann. Inst. H. Poincar\'e Probab. Statist. \textbf{33} (1997), 697--725, doi:10.1016/S0246-0203(97)80110-0.
\bibitem{BMT2001} V. S. Borkar, S. K. Mitter and S. Tatikonda, \emph{Optimal sequential vector quantization of Markov sources}, SIAM J. Control Optim. \textbf{40} (2001), 135--148, doi:10.1137/S0363012999365261.
\bibitem{Boyd2004} S. Boyd and L. Vandenberghe, \emph{Convex Optimization}, Cambridge University Press, 2004, doi:10.1017/CBO9780511804441.
\bibitem{DKY2024} Y. E. Demirci, A. D. Kara and S. Y\"uksel, \emph{Average cost optimality of partially observed MDPs: contraction of nonlinear filters and existence of optimal solutions and approximations}, SIAM J. Control Optim. \textbf{62} (2024), 2859--2883, doi:10.1137/24M1643736.
\bibitem{DGM2026} S. D\'iaz-Infante Velasco, Y. H. Garcia and J. A. Minj\'arez-Sosa, \emph{Recursive filtering and stochastic control under finite partition-based observations}, preprint (2026), arXiv:2608.13825.
\bibitem{Dobrushin1956} R. L. Dobrushin, \emph{Central limit theorem for nonstationary Markov chains. I}, Theory Probab. Appl. \textbf{1} (1956), 65--80, doi:10.1137/1101006.
\bibitem{Foster1985} J. Foster, R. M. Gray and M. O. Dunham, \emph{Finite-state vector quantization for waveform coding}, IEEE Trans. Inform. Theory \textbf{31} (1985), 348--359, doi:10.1109/TIT.1985.1057035.
\bibitem{Gneiting2007} T. Gneiting and A. E. Raftery, \emph{Strictly proper scoring rules, prediction, and estimation}, J. Amer. Statist. Assoc. \textbf{102} (2007), 359--378, doi:10.1198/016214506000001437.
\bibitem{KY2022} A. D. Kara and S. Y\"uksel, \emph{Near optimality of finite memory feedback policies in partially observed Markov decision processes}, J. Mach. Learn. Res. \textbf{23} (2022), no. 11, 1--46, arXiv:2010.07452.
\bibitem{Kim2026} M. Kim, \emph{Finite memory belief approximation for optimal control in partially observable Markov decision processes}, preprint (2026), arXiv:2601.03132.
\bibitem{Pages2015} G. Pag\`es and A. Sagna, \emph{Recursive marginal quantization of the Euler scheme of a diffusion process}, Appl. Math. Finance \textbf{22} (2015), 463--498, doi:10.1080/1350486X.2015.1091741.
\bibitem{PB2000} P. Poupart and C. Boutilier, \emph{Value-directed belief state approximation for POMDPs}, Proceedings of UAI (2000), 497--506, arXiv:1301.3887 (archived version).
\bibitem{PB2001} P. Poupart and C. Boutilier, \emph{Vector-space analysis of belief-state approximation for POMDPs}, Proceedings of UAI (2001), 445--452, arXiv:1301.2304 (archived version).
\bibitem{PB2002} P. Poupart and C. Boutilier, \emph{Value-directed compression of POMDPs}, Advances in Neural Information Processing Systems \textbf{15} (2002), 1547--1554.
\bibitem{GTFControl} Q. Qi, \emph{General Theta Foundations I: Acquired Geometry and Causal Resource Transfer}, continuation-curvature manuscript (2026), retained technical Companion V.
\bibitem{GTFAllocation} Q. Qi, \emph{General Theta Foundations I: Acquired Geometry and Causal Resource Transfer}, causal-allocation manuscript (2026), retained technical Companion W.
\bibitem{GTFBayes} Q. Qi, \emph{General Theta Foundations I: Acquired Geometry and Causal Resource Transfer}, Bayesian-curvature manuscript (2026), retained technical Companion H.
\bibitem{Tanaka2017} T. Tanaka, K.-K. K. Kim, P. A. Parrilo and S. K. Mitter, \emph{Semidefinite programming approach to Gaussian sequential rate-distortion trade-offs}, IEEE Trans. Automat. Control \textbf{62} (2017), 1896--1910, doi:10.1109/TAC.2016.2601148.
\bibitem{Varshney2011} K. R. Varshney, \emph{Bayes risk error is a Bregman divergence}, IEEE Trans. Signal Process. \textbf{59} (2011), 4470--4472, doi:10.1109/TSP.2011.2159500.
\end{thebibliography}

\end{document}
